\section{Combining meeting effects with a smooth curve}\label{sec:splice}
A piecewise-linear curve describes differences between meeting intervals. A smooth component can describe the broader shape across maturities: a level, a slope, or decaying exponential terms. We now add these two components and ask whether they can be driven by correlated shocks while preserving their shapes.

The source of the restriction is already visible in \eqref{eq:diffdrift}. If the total volatility is the sum of front-end and block volatilities, multiplying it by its maturity integral produces four terms: two from each component on its own and two cross terms. Those cross terms depend on covariance. They must be expressible using the functions available to the combined curve.

\subsection{The two curve components}
Use the piecewise-linear front end from \eqref{eq:front}, with levels $L_m(t)$ and finite-variation slopes $C_m(t)$. Let $x=T-t$ be time to maturity. The smooth block is a function $F(x,Z_t)$ of $x$ and a vector of stochastic coefficients $Z_t$. On maturity interval $I_m$, the total curve is
\begin{equation}
f(t,T)=L_m(t)+C_m(t)(T-T_m)+F(T-t,Z_t).\label{eq:splice}
\end{equation}

We take $F$ to be a polynomial plus a matrix-exponential component. Fix a polynomial degree $d\ge0$ and a matrix size $d_B\ge1$. The polynomial coefficients are $z_{0,0},\ldots,z_{0,d}$, and the exponential coefficients form a column $\zeta\in\R^{d_B}$. For the combined coefficient vector $z=(z_{0,0},\ldots,z_{0,d},\zeta)$, set
\begin{equation}
F(x,z)=\sum_{j=0}^d z_{0,j}x^j+\mathcal C e^{\mathcal A x}\zeta,
\qquad \mathcal A\in\R^{d_B\times d_B},\quad \mathcal C\in\R^{1\times d_B}.
\label{eq:block}
\end{equation}
For example, $d_B=1$, $\mathcal A=-\kappa$, and $\mathcal C=1$ give a single term $\zeta e^{-\kappa x}$. Matrix representations also include terms such as $x e^{-\kappa x}$. The exponential component may be omitted entirely.

We assume $\mathcal A$ is invertible, so the polynomial terms are kept separate from the nonzero-exponent terms. We also assume $(\mathcal C,\mathcal A)$ is observable: if $\mathcal C e^{\mathcal A x}v=0$ for every $x$, then $v=0$. This says that every retained coefficient direction changes the curve. Redundant coordinates can be removed before applying the criterion.

\subsection{Noise, covariance, and the drift discrepancy}
Use one common Brownian vector $W$ to describe both components. On interval $I_m$, let the row vector $V_m$ be the Brownian sensitivity of $L_m$. The front-end volatility at maturity $T\in I_m$ is therefore $V_m$. For each block coordinate, indexed by $k$, write its dynamics as
\[
\mathrm dZ^k_t=\mu_k(t)\dd t+\nu^k(t)\dd W_t.
\]
Here $\mu_k$ is its drift and $\nu^k$ is its row of Brownian sensitivities. A common driver represents arbitrary instantaneous correlations by inner products of these rows. Let $\nu^{0,j}$ denote the row for polynomial coefficient $z_{0,j}$, and let $\nu^\zeta$ be the matrix of rows for the exponential coefficients.

Two covariance quantities will enter the result. Define
\[
G_m=V_m-V_{m-1},\qquad \mathsf Q_{kl}=\nu^k(\nu^l)^\top.
\]
The vector $G_m$ is the change in front-end volatility as maturity crosses meeting $T_m$. The matrix $\mathsf Q$ is the instantaneous covariance of the block coefficients. For example, $\nu^kG_m^\top$ is the instantaneous covariance of block coordinate $k$ with the change in front-end noise at that meeting. If several factors contribute to the front end, $V_m$ includes all of them, so $G_m$ is their aggregate change.

To write the drift equation compactly, let $\varphi_k(x)$ be the function multiplying coordinate $z_k$ in \eqref{eq:block}: it is $x^j$ for a polynomial coordinate and the corresponding entry of $\mathcal C e^{\mathcal A x}$ for an exponential coordinate. Define its integral $\Phi_k(x)=\int_0^x\varphi_k(y)\dd y$. Then the block volatility is $\sum_k\varphi_k(x)\nu^k$. Its actual forward drift is $\sum_k \mu_k\varphi_k(x)-\partial_xF(x,Z)$, where the minus derivative appears because $x=T-t$ decreases as time passes at fixed $T$.

Subtracting the HJM drift generated by the block's own volatility gives
\begin{equation}
R(x)=\sum_k \mu_k\varphi_k(x)-\partial_xF(x,Z)
       -\sum_{k,l}\mathsf Q_{kl}\varphi_k(x)\Phi_l(x).\label{eq:residual}
\end{equation}
We call $R$ the block residual. It measures the discrepancy between the block's actual drift and the drift its volatility requires. For a model consisting of the block alone, consistency requires $R=0$. An independent piecewise-linear front end can absorb an affine residual by changing its level and slope drifts. Here and below, affine means a function of the form a constant plus a constant times $x$. For correlated components there is an additional covariance adjustment, which we now specify. Arguments $(t,\omega)$ of the coefficients and residual are suppressed in these pointwise formulas.

\subsection{The combined level noise and the criterion}
On the current maturity interval, the front-end noise $V_{j(t)}$ affects rates by a constant in maturity, just as the block level noise $\nu^{0,0}$ does. Their combined level noise is their sum. Define adjusted rows by
\begin{equation}
\nu^{0,0}_\sharp=\nu^{0,0}+V_{j(t)},\qquad \nu^k_\sharp=\nu^k\quad(k\ne(0,0)).\label{eq:effective}
\end{equation}
Let $\mathsf Q_{\sharp,kl}=\nu^k_\sharp(\nu^l_\sharp)^\top$, and let $R_\sharp$ be \eqref{eq:residual} with $\mathsf Q_\sharp$ in place of $\mathsf Q$. The block state and its drifts stay fixed in this calculation. The adjustment accounts for covariance with the current front-end level when testing the shape of the residual.

\subsection{Deriving the front-end drifts}
Expanding the drift equation identifies the front-end drifts that the criterion must make possible. Denote front-end and block volatilities by $\sigma^S(t,T)=V_m$ on $I_m$ and $\sigma^B(t,T)=\sum_k\varphi_k(T-t)\nu^k$. Define their maturity integrals by $\Sigma^S(t,T)=\int_t^T\sigma^S(t,u)\dd u$ and $\Sigma^B(t,T)=\int_t^T\sigma^B(t,u)\dd u$. Since the front-end volatility is constant on each interval, its integral has the form $TV_m+\chi_m$, with
\[
\Sigma^S(t,T)=TV_m+\chi_m,\qquad
\chi_{j(t)}=-tV_{j(t)},\qquad \chi_m-\chi_{m-1}=-T_mG_m.
\]
Multiplying total volatility by its integral gives the two individual-component terms and the following cross term on $I_m$:
\[
\mathcal K_m(t,T)=\sigma^B(t,T)\cdot\Sigma^S(t,T)+V_m\cdot\Sigma^B(t,T).
\]
Let $\mu_m^S(t,T)$ be the drift of $L_m(t)+C_m(t)(T-T_m)$ at fixed $T$. Subtract the actual block drift from the required total drift to obtain
\begin{equation}
\mu_m^S(t,T)=-R(T-t)+\mathcal K_m(t,T)+V_m\cdot\Sigma^S(t,T).
\label{eq:frontmatch}
\end{equation}

The right side consists of the negative block residual, the covariance cross term, and the front end's own required drift. It can be implemented using $L_m,C_m$ exactly when it is affine in maturity on each interval. The next theorem characterizes precisely when this happens.


The following theorem is a way to complete a model. Start with an existing block process $Z$ and proposed front-end noises $V_m$. Test the two stated conditions. If they hold, the theorem supplies front-end drifts that make the total curve consistent.

\begin{theorem}[Correlated splice criterion]\label{thm:splice}
For the block \eqref{eq:block}, assume $\mathcal A$ is invertible and $(\mathcal C,\mathcal A)$ is observable. Require the front-end slopes $C_m$ to have finite variation and no Brownian noise. Suppose the noises $V_m,\nu^k$ are predictable and square-integrable on each finite time horizon almost surely, the block drifts are progressively measurable and locally integrable, and $Z$ is a continuous semimartingale with those coefficients. There exist locally integrable front-end drifts for which \eqref{eq:splice} satisfies the integrated drift condition if and only if, for $dt\otimes dQ$-almost every $(t,\omega)$, both following conditions hold:
\begin{align}
\nu^{0,j}G_m^\top&=0\quad(1\le j\le d),&
\nu^\zeta G_m^\top&=0 \quad\text{for every }m\text{ with }T_m>t,\label{eq:steporth}\\
x&\longmapsto R_\sharp(x)\text{ is affine}
&&\text{on }[0,H-t].\label{eq:resaffine}
\end{align}
There is no step restriction on $\nu^{0,0}G_m^\top$.

The required front-end drifts are given by \eqref{eq:frontmatch}. If those drifts have already been prescribed, they must agree with that formula.
\end{theorem}

Condition \eqref{eq:steporth} tests what happens across each meeting. It requires zero covariance between a change in front-end volatility and each nonconstant block coordinate. The level coordinate is exempt because its contribution to the cross term is affine and can be absorbed by front-end drift.

Condition \eqref{eq:resaffine} tests the remaining smooth shape on the maturity domain. Since $R_\sharp$ is real analytic, affinity on any interval with nonempty interior extends throughout its analytic domain. For example, consider the independent block $F(x,z)=z_0+z_1x$ with Brownian noise of size $\eta_s$ in $z_1$ and none in $z_0$. Its volatility is $\eta_s x$, so the HJM drift it generates is $\eta_s^2x^3/2$. Its available parameter drifts have degree at most one. The residual therefore contains $-\eta_s^2x^3/2$ and cannot be affine unless $\eta_s=0$. Independence makes the first condition automatic in this example, but the second still rules out a stochastic slope.

A positive example is a block with only a level. All its covariance contributions are affine, so arbitrary covariance with the front-end levels can be accommodated by the chosen drifts. The difference between these examples explains the level exception in the theorem.

\begin{proof}[Proof of Theorem~\ref{thm:splice}]
First suppose the full equation holds. By imposing it at rational maturities and extending within each interval, work outside one time--probability null set. All remaining calculations are pointwise. Each interval expression for $\mathcal K_m$ is an analytic function of maturity that can be evaluated outside its own interval. Subtracting the expressions for adjacent intervals gives
\begin{equation}
\mathcal K_m-\mathcal K_{m-1}=\sum_k(\nu^kG_m^\top)
 \{\varphi_k(x)(T-T_m)+\Phi_k(x)\},\qquad x=T-t.\label{eq:crossjump}
\end{equation}
The right side must be affine. Its nonzero-exponent part is
\[
\mathcal C e^{\mathcal A x}\{(x+t-T_m)I+\mathcal A^{-1}\}\nu^\zeta G_m^\top.
\]
An exponential polynomial with nonzero exponents cannot equal a polynomial unless its nonzero-exponent part vanishes. Observability then forces $\nu^\zeta G_m^\top=0$; the elementary identity needed for this last step is proved in Appendix~\ref{app:proofs}. In the remaining polynomial part, the term indexed by $j$ has degree $j+1$ and leading coefficient $(j+2)\nu^{0,j}G_m^\top/(j+1)$. The coefficient of $x^{d+1}$ must first vanish, so $\nu^{0,d}G_m^\top=0$. Removing that term leaves the same argument at degree $d$, and continuing down to degree two proves \eqref{eq:steporth}. The level term is affine and survives.

Put $\bar\omega_k=\nu^kV_{j(t)}^\top$. Under \eqref{eq:steporth}, the non-level part of $\mathcal K_m$ is the same on every later interval and equals
\begin{equation}
\mathcal X(x)=\sum_{j=1}^d\bar\omega_j\frac{j+2}{j+1}x^{j+1}
 +\mathcal C\{x e^{\mathcal A x}+\mathcal A^{-1}(e^{\mathcal A x}-I)\}\bar\omega_\zeta.
\label{eq:crosscommon}
\end{equation}
The level part of $\mathcal K_m$ is affine, with adjacent-interval difference
$(\nu^{0,0}G_m^\top)(2T-t-T_m)$. Changing the level covariance as in \eqref{eq:effective} gives the exact identity
\begin{equation}
R-R_\sharp=\mathcal X+(2\bar\omega_0+|V_{j(t)}|^2)x.\label{eq:effectiveidentity}
\end{equation}
On the current maturity interval, \eqref{eq:frontmatch} consequently forces $R_\sharp$ to be affine. It is real analytic in $x$, so it is affine throughout its domain. This proves necessity.

Conversely, \eqref{eq:steporth} and \eqref{eq:resaffine}, together with \eqref{eq:effectiveidentity}, make the right side of \eqref{eq:frontmatch} affine on every interval. Choose the drift of $L_m$ to be its value at $T_m$ and the drift of $C_m$ to be its slope, using the interval's analytic expression even when $T_m<t$. This gives the differentiated HJM equation and hence its integrated form.

It remains to check integrability, since coefficients were not assumed bounded. On a finite horizon the coordinate functions and their derivatives are bounded on the maturity range. The chosen drifts are finite sums of the block drifts, state coordinates, and products of pairs of noise rows with bounded deterministic multipliers. Continuity of $Z$ bounds its paths on compact time intervals. Each noise product is integrable by
$|v\cdot w|\le(|v|^2+|w|^2)/2$. The front-end drift integrals and noise integrals therefore define $L_m,C_m$ with the stipulated regularity. This also provides bounds on the maturity suprema needed to integrate the drift identity. 
\end{proof}

The integrability hypothesis allows stochastic volatility scales. For example, if $v_t$ is an existing continuous nonnegative variance process, multiplying bounded noise rows by $\sqrt{v_t}$ gives square-integrable rows on each finite horizon, path by path. The theorem can then be applied with those rows. Existence of the block and variance processes is an input; the construction supplies the front end.

\subsection{A stronger restriction for sums of pure exponentials}
If $\mathcal A$ is diagonalizable, its coordinates can be chosen so that the exponential part is a sum of terms $e^{\upsilon x}$ with nonzero exponents, without additional powers of $x$. Covariance between these terms and a polynomial coordinate produces terms such as $x e^{\upsilon x}$. Since those shapes are absent from the drift basis, their coefficients must vanish. This yields a further consequence of the residual condition.

\begin{corollary}\label{cor:diagonal}
Under Theorem~\ref{thm:splice}, if $\mathcal A$ is diagonalizable over $\mathbb C$, consistency implies
\begin{equation}
\nu^\zeta(\nu^{0,j})^\top=0\quad(j\ge1),\qquad
\nu^\zeta(\nu^{0,0}+V_{j(t)})^\top=0.\label{eq:effectiveorth}
\end{equation}
\end{corollary}
Thus the exponential noise must have zero covariance with every positive-degree polynomial coordinate and with the \emph{combined} level noise. It need not be separately orthogonal to the block level and the current front-end level if their covariances cancel.

\begin{proof}
In an observable complex eigenbasis, the nonzero-exponent coordinate functions are $\mathcal C_i e^{\upsilon_i x}$ with $\mathcal C_i\ne0$. Terms $x^ke^{\upsilon_i x}$, $k\ge1$, in $R_\sharp$ can only arise from polynomial--exponential covariance. The coefficient is
$-\mathcal C_i\{\mathsf Q_{\sharp,(k,i)}/\upsilon_i+\mathsf Q_{\sharp,(k-1,i)}/k\}$, with absent indices interpreted as zero. At $k=d+1$ the first entry is absent, so the coefficient forces the degree-$d$ covariance to vanish. At $k=d$ that entry is now zero, forcing the degree-$(d-1)$ covariance to vanish. Repeating through $k=1$ also removes the level covariance. Products of two pure exponentials can change exponents but have no positive polynomial power, so resonances do not alter this argument.
\end{proof}

If no block level is present, append one with zero state, drift, and noise. Then the final restriction in \eqref{eq:effectiveorth} becomes $\nu^\zeta V_{j(t)}^\top=0$. In that case the exponential block must be uncorrelated with the current front-end noise as well as with future changes in that noise.

If two front factors have cancelling maturity steps along the same driver, their individual covariations need not vanish: their aggregate $G_m$ is zero. If the exponential representation is not observable, quotient out its unobservable subspace and apply the theorem to the resulting noise rows. For blocks containing terms such as $xe^{-\beta x}$, those functions are already available to the drift. The full residual test in Theorem~\ref{thm:splice} then applies, rather than the stronger conclusion for pure exponentials.

Section~\ref{app:blocks} relates this residual formulation to classical consistency calculations. The examples there show how absorbing an affine residual changes classical consistency conditions while retaining the positive-semidefinite covariance requirement.

\subsection{Identifiability of the component levels and slopes}\label{sec:normalization}
The decomposition of the total curve is not unique. For any common level shift $u(t)$, replacing $L_m$ by $L_m+u$ for every $m$ and $z_{0,0}$ by $z_{0,0}-u$ leaves $f$ unchanged. Taking $u=z_{0,0}$ gives a normalization with no block level. Its Brownian row becomes $V_m+\nu^{0,0}$, while $G_m$ and the combined level noise are unchanged. This explains why the criterion uses the combined noise: that quantity belongs to the curve, rather than to a particular allocation of its level between components.

At the level of curve functions, the linear block term is also redundant when $d\ge1$. To set $z_{0,1}$ to zero, replace
\[
C_m\mapsto C_m+z_{0,1},\qquad
L_m\mapsto L_m+z_{0,1}(T_m-t).
\]
For a dynamic model this normalization preserves our finite-variation slope convention only if $z_{0,1}$ itself has finite variation. A Brownian block slope cannot be transferred to $C_m$ within the stated model class. Thus one can always normalize the block level to zero; a zero block-slope normalization requires that extra restriction. Calibration should fix a convention for these overlapping coefficients before interpreting them individually.
